\chapter{Outlook and Open Problems}
\label{chap:outlook}

The preceding chapters developed a progression of temporal-task planning and
coordination frameworks, starting from the single-robot automata-based
synthesis reviewed in Chapter~\ref{chap:preliminaries}, extending to
bottom-up and top-down multi-robot coordination in
Chapters~\ref{chap:bottom-up} and~\ref{chap:top-down}, and then addressing
partial knowledge, limited communication, and dynamic uncertainty in
Chapters~\ref{chap:communication} and~\ref{chap:uncertainty}. Together,
these developments establish a methodological foundation for formal
multi-robot coordination, while also exposing limits that cannot be resolved
by incremental algorithmic improvements alone. As task structures, robot teams,
and deployment settings become more complex, the central questions shift from
nominal satisfaction on fixed abstractions to scalability, continual adaptation, semantic
incompleteness, data-driven modeling, and system-level validation.
This chapter organizes these questions by their distinct sources of
difficulty. Section~\ref{sec:ch07_research_challenges} groups five research
challenges and open problems. Section~\ref{sec:ch07_scalability} focuses on
computational and structural scale: when and how temporal-task problems can be
decomposed without losing the dependencies that make the specification
meaningful. Section~\ref{sec:ch07_continual} considers nonstationary
execution, where the modeling framework remains fixed but tasks, targets,
resources, or communication conditions change over time.
Section~\ref{sec:ch07_open_world} addresses a different issue: partial
knowledge and open-world semantics, in which the symbolic model itself may be
incomplete or extensible. Section~\ref{sec:ch07_learning} then discusses
how learned components may be introduced into this formal pipeline.
Section~\ref{sec:ch07_deployment} turns to runtime assurance, human
supervision, and benchmark design. Section~\ref{sec:ch07_concluding}
concludes the chapter.

%=================================================
\section{Research challenges and open problems}
\label{sec:ch07_research_challenges}

\subsection{Scalability and compositional interfaces}
\label{sec:ch07_scalability}

A recurring theme throughout this monograph is that the expressive strength of
temporal-task specifications comes with a substantial computational burden.
Product constructions, joint transition systems, and centralized assignment
models remain conceptually clean, but their size grows quickly with the number
of robots, task clauses, workspace regions, and synchronization requirements.
Chapters~\ref{chap:bottom-up} and~\ref{chap:top-down} show that distribution
alone does not remove this difficulty: the coupling reappears through local
subtasks, shared resources, ordering constraints, and heterogeneous
capabilities. The open problem is therefore not simply to design faster
solvers, but to understand which formal structures admit compact
compositional representations.

(I) \emph{Structural limits of monolithic synthesis.}
A first unresolved issue is to characterize when monolithic synthesis is
fundamentally unavoidable. Some tasks contain sparse temporal coupling, so
that only a few robots or subtasks interact at any time. Others appear large
but have weak effective coupling, because most interactions are irrelevant to
feasibility or quality. In contrast, tightly synchronized specifications may
force centralized reasoning regardless of implementation. A mature
scalability theory should identify these regimes through structural criteria,
not only through worst-case complexity bounds. Such criteria would clarify
when decomposition is sound, when it is merely heuristic, and when it cannot
preserve the intended task semantics.

(II) \emph{Compositional contracts for tasks and teams.}
A second difficulty is the certification of interfaces between subproblems.
When a global specification is decomposed into local tasks, the planner must
specify what each local solution promises to the rest of the system and what
assumptions it requires in return. These contracts may involve temporal
ordering, resource availability, communication events, safety margins, or
capability constraints. The challenge becomes sharper for heterogeneous
teams, where robots differ in sensing, mobility, endurance, and access to
task-relevant actions. A general notion of compositional soundness remains
missing across symbolic decomposition, task assignment, communication-aware
coordination, and continuous trajectory coupling.

(III) \emph{Cross-layer abstraction and refinement.}
A related challenge concerns the interaction among planning layers. Discrete
task planning, motion planning, communication planning, and resource
allocation are usually modeled at different resolutions and time scales. If
these layers are connected only through case-specific interfaces, local refinements
may invalidate decisions made elsewhere. Future methods should therefore provide
abstractions that are both compact and refinable: they should expose enough
information to support symbolic reasoning, while retaining the geometric,
resource, and communication constraints needed for execution. The long-term
goal is a cross-layer synthesis architecture in which local refinements do not
force complete recomputation of the full coordination policy.

%=================================================
\subsection{Continual coordination under nonstationarity}
\label{sec:ch07_continual}

Scalability concerns the size and structure of the planning problem. A
separate difficulty arises when the planning problem changes during execution.
Here the relevant model class may still be known, but targets move,
communication conditions vary, resources are consumed, risks change, or new
requests arrive while the team is already committed to ongoing actions.
Several chapters point in this direction, especially through online task
updates, intermittent communication, and dynamic constraints. The open problem
is to treat revision as part of the formal problem rather than as an external
repair step.

(I) \emph{Persistent obligations and online task arrival.}
Many missions do not end after a single specification has been satisfied. New
subtasks may arrive before earlier commitments are complete, priorities may be
revised, and interrupting an ongoing action may be costly or unsafe.
Coordination must therefore reason over completed, active, and pending
obligations simultaneously. This calls for specification models that can
represent task queues, persistent mission context, and controlled priority
revision. Feasibility should then mean more than eventual satisfaction of a
new request; it should also account for whether the request can be integrated
without invalidating commitments already in force.

(II) \emph{Local repair, reassignment, and global revision.}
A second open question is how to choose the appropriate scale of reaction.
Some changes can be handled by local trajectory repair, others require task
reassignment, and still others invalidate the current symbolic plan. Treating
all changes by full replanning is usually too expensive, while treating all
changes locally can miss global conflicts. Future coordination methods need
criteria that distinguish these cases. Such criteria should depend on the
location of the change, the temporal role of the affected task, the available
team redundancy, and the coupling between the affected component and the remaining
mission.

(III) \emph{Guarantees over adaptation sequences.}
Repeated adaptation changes the meaning of performance and correctness. A
sequence of individually feasible repairs may still produce oscillatory task
assignments, excessive disruption, or poor long-horizon mission quality. The
analysis should therefore move from isolated replanning events to adaptation
sequences. Relevant properties include bounded disruption to ongoing
execution, monotonic or recoverable logical progress, stability of team
assignments, and cumulative performance over extended horizons. These
properties remain underdeveloped for temporal-task planning, especially when
formal task logic is combined with resource, communication, and motion
constraints.

%=================================================
\subsection{Open-world task semantics}
\label{sec:ch07_open_world}

Nonstationarity concerns changes within an identifiable modeling framework.
Open-world planning concerns a deeper form of incompleteness: the relevant
states, propositions, objects, or constraints may not be fully known when
planning begins. This distinction is important. A nonstationary problem asks
how to react when known quantities change; an open-world problem asks how to
reason when the symbolic description itself is provisional. In realistic
multi-robot missions, maps may be incomplete, task-relevant objects may be
discovered during execution, and the meaning of a proposition may depend on
context or newly acquired information.

(I) \emph{Provisional symbolic models.}
Many formal planning frameworks assume that the proposition set and state
abstraction are fixed before synthesis. In exploration, inspection, or search
missions, this assumption is often too strong. The team may discover new
regions, revise traversability, or add task-relevant entities after execution
has started. The planner therefore needs symbolic models that can be expanded
and corrected while retaining a traceable relation to earlier decisions. This
requires mechanisms for recording which commitments were made under which
model assumptions, and for identifying when a model update preserves,
weakens, or invalidates those commitments.

(II) \emph{Semantic status of task propositions.}
A proposition in temporal logic is often treated as an atomic truth value, but
in robotic deployment it may encode a semantic judgment: a region is safe, an
object is relevant, a victim has been found, or a request has been fulfilled.
Such judgments may depend on context, evidence quality, operational
convention, or information shared across the team. Open-world semantics should
therefore make the status of propositions explicit, for example through
confidence, provenance, validity conditions, or levels of commitment. The
central question is what temporal-task satisfaction should mean when the task
vocabulary is not simply observed, but constructed and revised during
execution.

(III) \emph{Controlled expansion of task models.}
The broad open problem is to formulate temporal-task planning when the model
boundary itself can grow. New objects, propositions, agents, and constraints
may enter the mission after planning has begun. Unrestricted expansion would
make formal reasoning unstable; forbidding expansion would make the framework
unrealistic. A useful theory must therefore specify how model growth is
controlled: which additions are conservative, which require revalidation of
past decisions, and which force a change of task interpretation. This problem
is distinct from learning or runtime verification; it concerns the semantics
of planning under incomplete and extensible symbolic descriptions.

%=================================================
\subsection{Learning-augmented planning with formal guarantees}
\label{sec:ch07_learning}

Learning-based components are attractive because many quantities needed for
deployment cannot be modeled accurately in advance. Perception may
supply task labels, predictors may estimate target motion or risk evolution,
and data-driven heuristics may accelerate assignment or motion planning. The
open problem is not whether learning can improve temporal-task planning, but
how learned quantities should enter a formal architecture without obscuring
which assumptions support each decision.

(I) \emph{Roles of learning in the planning pipeline.}
Learning can appear at several layers, and these roles should not be
conflated. A learned perception module affects the truth evaluation of
propositions. A learned dynamics or risk model affects prediction. A learned
abstraction or heuristic affects computational search. A learned motion
primitive affects execution feasibility. Each role has different consequences
for formal reasoning. Future frameworks should classify learned components by
where they enter the pipeline, whether they alter the symbolic model or only
guide search, and how their errors propagate to task-level decisions.

(II) \emph{Interfaces between learned and symbolic components.}
The central integration problem is the design of interfaces. A planner should
not treat learned outputs as exact symbolic facts unless their reliability is
specified. Useful interfaces may include calibrated confidence,
validity domains, uncertainty sets, abstention mechanisms, and recovery
policies when learned outputs become unreliable. These interfaces should be
modular, so that the formal planner can reason about the effect of a learned
component without requiring a complete internal model of that component. This
is especially important when multiple learned modules influence different
parts of the same temporal-task specification.

(III) \emph{Guarantee notions for symbolic and probabilistic systems.}
Classical guarantees are usually stated with respect to a fixed abstraction.
Learning introduces statistical error, distribution shift, and context
dependence, so the guarantee language itself may need to change. Possible
forms include guarantees conditioned on calibrated confidence, bounds over
specified mismatch classes, probabilistic task satisfaction under verified
assumptions, or safety envelopes that restrict when learned information may be
used. The unresolved question is whether hybrid planning should approximate
classical formal guarantees as closely as possible, or develop guarantee
notions tailored to systems whose models are partly inferred from data.

%=================================================
\subsection{Runtime assurance and deployment evaluation}
\label{sec:ch07_deployment}

The final set of open problems arises at the level of deployed systems.
Planning algorithms may be correct relative to their models and still fail to
support dependable operation when sensing, communication, human input, and
execution constraints interact. This subsection therefore focuses on evaluation
and assurance during operation, rather than on the internal semantics of model
incompleteness or the statistical properties of learned modules.

(I) \emph{Runtime assurance and assumption monitoring.}
Offline verification should be complemented by mechanisms that monitor the
assumptions under which a plan remains valid. These assumptions may concern
communication availability, proposition evaluation, localization accuracy,
resource margins, or the validity domain of a learned component. Runtime
assurance should detect violations, estimate their effect on task progress,
and select among continuation, repair, fallback, or suspension. This shifts
verification from a single pre-execution certificate to an operational process
that tracks whether the deployed system remains inside the conditions for
which its formal claims are meaningful.

(II) \emph{Human-supervised temporal-task execution.}
Human operators are often continuing participants rather than only task
designers. They may add requests, change priorities, override autonomous
decisions, or ask for explanations during execution. Future temporal-task
frameworks should model supervisory authority explicitly: which decisions can
be overridden, how intervention affects existing commitments, and how the
planner explains the consequences of a requested change. The goal is not to
replace formal consistency with unstructured human control, but to make human
supervision a structured part of the coordination problem.

(III) \emph{Benchmarks and deployment metrics.}
Progress in this area also requires better evaluation practice. Correctness,
completion rate, and computation time remain useful metrics, but they do not
capture long-horizon behavior under uncertainty, communication limits,
revision, and supervision. Benchmark suites should therefore measure
coordination stability, adaptation cost, communication burden, recoverability,
interpretability, and cumulative mission quality. They should also make it
possible to compare symbolic, optimization-based, and learning-augmented
methods on shared scenarios rather than on isolated demonstrations.

%==============================
\section{Concluding perspective}
\label{sec:ch07_concluding}

The open problems above are connected, but they are not interchangeable.
Scalability concerns structural decomposition; continual coordination concerns
change over time; open-world semantics concerns incomplete symbolic models;
learning concerns data-driven components inside the planning pipeline; and
deployment evaluation concerns assurance in the executed system. Keeping
these distinctions clear is essential for future progress.
The long-term objective is to extend the discipline of formal temporal-task
planning without assuming away the conditions that make robotic deployment
difficult. Future frameworks will likely combine compositional synthesis,
adaptive abstraction, learned inference, runtime assurance, and human
supervision. Their success should be measured not only by whether temporal
specifications can be satisfied efficiently, but also by whether correctness,
robustness, and interpretability remain meaningful as the world, the model,
and the coordination structure evolve during execution.
